\capitulo{8}{Lineas de trabajo futuras}

Como posibles líneas de trabajo futuras se plantean las siguientes:
\begin{enumerate}
    \item \textbf{Ampliar el número de enzimas consideradas.} Incorporando otras vías metabólicas además de las seis isoenzimas CYP450 utilizadas en este trabajo, con el objetivo de aumentar la cobertura de medicamentos analizados y evitar posibles malas clasificaciones por parte del algoritmo.
    \item \textbf{Mejorar el algoritmo de clasificación del riesgo.} Incorporando información cuantitativa sobre la intensidad de la inhibición, inducción o afinidad de cada principio activo por las distintas enzimas metabólicas.
    \item \textbf{Mejorar la extracción de información farmacológica mediante técnicas de procesamiento de lenguaje natural.}Incorporando métodos de procesamiento de texto que permitan analizar automáticamente la información no estructurada disponible en bases de datos como DrugBank y en las fichas técnicas oficiales para identificar con mayor precisión el papel de cada principio activo como sustrato, inhibidor o inductor de las diferentes isoenzimas del sistema CYP450.
    \item \textbf{Integrar información farmacogenética.} Permitiendo adaptar el análisis a los polimorfismos genéticos del paciente y avanzar hacia una herramienta orientada a la medicina personalizada.
    \item \textbf{Actualizar automáticamente la base de conocimiento.} Implementando mecanismos que permitan sincronizar periódicamente la información con las nuevas versiones de DrugBank, CIMA y otras bases de datos farmacológicas.
    \item \textbf{Ampliar las fuentes de información.} Incorporando nuevas bases de datos especializadas sobre interacciones, farmacovigilancia y evidencia clínica para mejorar la robustez del sistema.
    \item \textbf{Validar la herramienta en un entorno clínico.} Comparando sus resultados con los obtenidos por herramientas ya existentes, utilizadas por el personal facultativo y con la valoración de profesionales sanitarios para evaluar su utilidad asistencial.
    \item \textbf{Optimizar el rendimiento y la escalabilidad de la aplicación.} Mediante técnicas de almacenamiento y procesamiento más eficientes que permitan manejar un mayor volumen de datos y consultas simultáneas.
    \item \textbf{Implementación de la herramienta con Streamlit}. Se plantea migrar la aplicación al framework Streamlit con el objetivo de facilitar su despliegue y permitir el acceso a todas las funcionalidades desarrolladas, evitando las limitaciones de tiempo de ejecución y recursos presentes en la plataforma Render.
\end{enumerate}
